% Enrichment + Inductive + Computable (Ch 15-17)
% This file is for agent editing — will be merged back

\chapter{The Semantic Enrichment Theorem}
\label{ch:enrichment}

The most powerful extension of CTTT: using the cubical path space
to \emph{enrich} the \v{C}ech complex, yielding a finer cohomology
that distinguishes programs the basic theory cannot.

\section{Enrichment of the Program Category}

Recall that the program category $\mathcal{P}$ has Python functions as
objects and equivalence witnesses as morphisms.  So far we have treated
$\Hom(f,g)$ as a \emph{set}: either empty (non-equivalent) or a
singleton (equivalent).  The enrichment lifts $\Hom(f,g)$ from
$\Set$ to a richer category that records \emph{how} $f$ and $g$ are
related.

\begin{definition}[Duck-Type Lattice]
\label{def:duck-lattice}
Let $\mathcal{D}$ be the lattice of duck types ordered by subtyping:
\[
  \mathcal{D} \coloneqq (\{\mathsf{Addable}, \mathsf{Iterable},
    \mathsf{Comparable}, \mathsf{Hashable}, \ldots\},\; \sqsubseteq)
\]
where $D_1 \sqsubseteq D_2$ iff every operation required by $D_1$ is
also required by $D_2$.  This is a finite lattice with meet
$D_1 \sqcap D_2$ (intersection of required operations) and join
$D_1 \sqcup D_2$ (union).
\end{definition}

\begin{definition}[Enriched Hom]
\label{def:enriched-hom}
The \emph{$\mathcal{D}$-enriched hom} between programs $f$ and $g$ is:
\[
  \underline{\Hom}(f,g) \coloneqq \bigsqcap_{U_i \in \covers}
    \{D \in \mathcal{D} \mid
    \Sem_f|_{U_i} =_D \Sem_g|_{U_i}\} \;\in\; \mathcal{D}
\]
That is, the meet (in the duck-type lattice) of all duck types under
which $f$ and $g$ agree on each fiber.  If $f$ and $g$ agree as
$\mathsf{Addable}$ on the $\mathsf{Int}$ fiber but only as
$\mathsf{Comparable}$ on the $\mathsf{Str}$ fiber, then
$\underline{\Hom}(f,g) = \mathsf{Comparable}$ (the coarser agreement).
\end{definition}

The ordinary (boolean) hom is recovered by forgetting the enrichment:
$\Hom(f,g) \neq \varnothing$ iff
$\underline{\Hom}(f,g) \neq \bot_{\mathcal{D}}$.

\section{Enriched Cohomology}

\begin{definition}[Path-Enriched Cochain]
An \emph{enriched 0-cochain} assigns to each fiber $U_i$ not just a
boolean (isomorphism or not) but a \emph{path type}:
\[
  C^0_\mathsf{enr}(\covers, \Iso) \coloneqq \prod_i
    \|\Path(\Sem_f|_{U_i}, \Sem_g|_{U_i})\|
\]
where $\|\cdot\|$ is propositional truncation (we care about
existence, not the specific path).
\end{definition}

The enrichment adds information: we know not just \emph{whether} two
programs agree on a fiber, but \emph{how} they agree (which axioms
connect them).

\begin{definition}[Enriched 1-Cochain]
The enriched 1-cochain group is:
\[
  C^1_\mathsf{enr}(\covers, \Iso) \coloneqq \prod_{i < j}
    \Path\bigl(\restr{p_i}{U_{ij}},\; \restr{p_j}{U_{ij}}\bigr)
\]
where $U_{ij} = U_i \cap U_j$ is the overlap fiber.  An element of
$C^1_\mathsf{enr}$ records, for each overlap, whether the local
equivalence paths are \emph{themselves} path-connected---i.e.\
whether the axioms used on adjacent fibers are compatible.
\end{definition}

\begin{definition}[Enriched Coboundary]
The enriched coboundary $\delta^0_\mathsf{enr} : C^0_\mathsf{enr}
\to C^1_\mathsf{enr}$ assigns to each overlap $U_{ij}$ the
\emph{transition path}:
\[
  (\delta^0_\mathsf{enr}\, \varphi)_{ij} = \restr{p_j}{U_{ij}}
    \cdot (\restr{p_i}{U_{ij}})^{-1}
\]
where $p_i$ and $p_j$ are the local paths on fibers $U_i$ and $U_j$.
A cocycle $\varphi \in \ker \delta^0_\mathsf{enr}$ is a family of
local paths whose transition maps are all trivial (the identity path).
\end{definition}

\begin{theorem}[Enriched Descent]
\label{thm:enriched-descent}
The enriched cohomology $\Hone_\mathsf{enr}$ is a refinement of the
ordinary $\Hone$:
\[
  \Hone_\mathsf{enr} = 0 \implies \Hone = 0
\]
but not vice versa.  The enriched version detects \emph{path
incompatibilities}: situations where local paths exist but use
incompatible axioms on overlaps.
\end{theorem}

\begin{proof}
There is a natural forgetful map
$\pi : C^\bullet_\mathsf{enr} \to C^\bullet$ given by truncating
each path type to its mere proposition (existence).  This map is a
cochain morphism ($\pi \circ \delta_\mathsf{enr} =
\delta \circ \pi$) and hence induces a map on cohomology:
\[
  \pi^* : \Hone_\mathsf{enr} \to \Hone.
\]
If $\Hone_\mathsf{enr} = 0$, then every enriched 1-cocycle is a
coboundary.  Applying $\pi^*$, every ordinary 1-cocycle is also a
coboundary, so $\Hone = 0$.

For the converse failure, consider two fibers $U_1, U_2$ with
local paths $p_1, p_2$ such that $p_1|_{U_{12}}$ uses axiom
$\alpha$ (e.g.\ commutativity of $+$) and $p_2|_{U_{12}}$ uses
axiom $\beta$ (e.g.\ associativity of $\cdot$) where
$\alpha$ and $\beta$ are not themselves connected by a path in the
axiom space.  Then $\Hone = 0$ (both fibers pass the boolean check)
but $\Hone_\mathsf{enr} \neq 0$ (the transition path
$p_2|_{U_{12}} \cdot (p_1|_{U_{12}})^{-1}$ is a non-trivial
loop in the axiom groupoid).
\end{proof}

\begin{example}[Path Incompatibility]
Program $f$ uses \texttt{isinstance(x, int)} to dispatch.  On the
$\mathsf{Int}$ fiber, the equivalence path uses the $\mathsf{add}$
commutativity axiom.  On the $\mathsf{Str}$ fiber, the path uses
string concatenation (which is NOT commutative).

The ordinary $\Hone$ says: both fibers are locally equivalent, so
$\Hone = 0$.  But the enriched $\Hone$ detects that the transition
path on the Int$\cap$Str overlap is inconsistent (commutativity on
one side, non-commutativity on the other).

This reveals a subtle bug: the program treats \texttt{+} as
commutative everywhere, but it's only commutative for integers.
\end{example}

\section{The Enrichment--Degeneration Tradeoff}

\begin{proposition}[Enrichment Hierarchy]
\label{prop:enrichment-hierarchy}
There is a tradeoff between enrichment and computability, forming
a three-level hierarchy:
\begin{enumerate}
  \item \textbf{Boolean enrichment} (ordinary $\Hone$): decidable,
    computable in $O(n^3)$ via GF(2) rank.  This is the
    $\mathcal{D} = \{0,1\}$ case.
  \item \textbf{Duck-type enrichment} ($\underline{\Hom}$ over
    $\mathcal{D}$): decidable when $\mathcal{D}$ is finite, since
    the meet computation is a finite lattice operation.
  \item \textbf{Path enrichment} (enriched $\Hone$): semi-decidable
    (each local path check is a Z3 query, which may timeout).
  \item \textbf{Full cubical enrichment} (tracking the specific
    path constructors used): undecidable in general, but provides
    the most information.
\end{enumerate}
\end{proposition}

In practice, we use boolean enrichment (fast) as a first pass, then
duck-type enrichment (still fast, more informative) to determine
\emph{at what level of abstraction} $f$ and $g$ agree, and finally
path enrichment (slower) on fibers where the duck-type check is
ambiguous.  This is the spectral sequence strategy applied to the
enrichment dimension.


\chapter{Inductive Families and Dependent Pattern Matching}
\label{ch:inductive-families}

The most ambitious extension: using \emph{inductive families} (dependent
inductive types indexed by values) to capture program invariants
that vary with the input.  We show that the \texttt{OTerm} type
from the implementation is itself an inductive family, and that the
normalizer's phase-by-phase case analysis is exactly the dependent
elimination principle.

\section{Inductive Families in CTTT}

\begin{definition}[Inductive Family]
An \emph{inductive family} $T : I \to \Type$ indexed by $I$ is a type
defined by constructors whose return type depends on the index:
\begin{align}
  \mathsf{nil} &: T(\mathsf{empty}) \\
  \mathsf{cons} &: \prod_{i : I} A \to T(i) \to T(\mathsf{extend}(i))
\end{align}
\end{definition}

\begin{example}[Length-Indexed Lists]
The type $\mathsf{Vec}(A, n)$ of lists of length exactly $n$:
\begin{align}
  \mathsf{nil} &: \mathsf{Vec}(A, 0) \\
  \mathsf{cons} &: A \to \mathsf{Vec}(A, n) \to \mathsf{Vec}(A, n+1)
\end{align}
A program that takes a $\mathsf{Vec}(A, n)$ and returns a
$\mathsf{Vec}(B, n)$ is \emph{guaranteed} to preserve length.
\end{example}

\section{OTerm as an Inductive Family}

The \texttt{OTerm} type in \texttt{denotation.py} is the
computational manifestation of an inductive family indexed by
\emph{arity} (number of free variables).

\begin{definition}[The OTerm Family]
\label{def:oterm-family}
Define the inductive family $\mathsf{OTerm} : \Nat \to \Type$
indexed by arity $n \in \Nat$ via the following constructors:
\begin{align}
  \mathsf{OVar}(k)       &: \mathsf{OTerm}(n)
    & \text{where } k < n \\
  \mathsf{OLit}(v)       &: \mathsf{OTerm}(0)
    & \text{(literals are closed)} \\
  \mathsf{OOp}(f, \vec{a}) &: \mathsf{OTerm}(n)
    & \text{where } \vec{a} : \mathsf{OTerm}(n)^m \\
  \mathsf{OFold}(f, z, c) &: \mathsf{OTerm}(n)
    & \text{where } z, c : \mathsf{OTerm}(n) \\
  \mathsf{OCase}(t, b_1, b_2) &: \mathsf{OTerm}(n)
    & \text{where } t, b_1, b_2 : \mathsf{OTerm}(n) \\
  \mathsf{OFix}(b)       &: \mathsf{OTerm}(n)
    & \text{where } b : \mathsf{OTerm}(n{+}1) \\
  \mathsf{OSeq}(\vec{e}) &: \mathsf{OTerm}(n)
    & \text{where } \vec{e} : \mathsf{OTerm}(n)^m \\
  \mathsf{OLam}(b)       &: \mathsf{OTerm}(n)
    & \text{where } b : \mathsf{OTerm}(n{+}k) \\
  \mathsf{OMap}(f, c)    &: \mathsf{OTerm}(n)
    & \text{where } f : \mathsf{OTerm}(n{+}1),\; c : \mathsf{OTerm}(n) \\
  \mathsf{OQuotient}(t, R)  &: \mathsf{OTerm}(n)
    & \text{where } t : \mathsf{OTerm}(n)
\end{align}
Note that $\mathsf{OFix}$ and $\mathsf{OLam}$ \emph{change the index}:
their body has arity $n{+}1$ (or $n{+}k$), reflecting the bound
variable(s).  This is what makes $\mathsf{OTerm}$ a genuine inductive
family, not merely a simple inductive type.
\end{definition}

\begin{remark}[Implementation Correspondence]
In \texttt{denotation.py}, the arity index is implicit: \texttt{OVar}
stores a name string rather than a de Bruijn index, and \texttt{OLam}
stores parameter names.  The \texttt{canon()} method performs
alpha-normalization (replacing parameter names with positional
indices \texttt{\_b0}, \texttt{\_b1}, \ldots), which recovers the
de Bruijn representation and thus the indexed family structure.
\end{remark}

\section{The Elimination Principle}

\begin{definition}[OTerm Eliminator]
\label{def:oterm-elim}
The elimination principle for $\mathsf{OTerm}$ is a dependent
recursor: for any type family $P : \prod_{n:\Nat}
\mathsf{OTerm}(n) \to \Type$, given methods
\begin{align}
  m_{\mathsf{Var}} &: \prod_{n:\Nat}\prod_{k < n} P(n, \mathsf{OVar}(k)) \\
  m_{\mathsf{Lit}} &: \prod_{v} P(0, \mathsf{OLit}(v)) \\
  m_{\mathsf{Op}} &: \prod_{n,f,\vec{a}}
    \bigl(\prod_i P(n, a_i)\bigr) \to P(n, \mathsf{OOp}(f,\vec{a})) \\
  m_{\mathsf{Fold}} &: \prod_{n,f,z,c}
    P(n,z) \to P(n,c) \to P(n, \mathsf{OFold}(f,z,c)) \\
  m_{\mathsf{Case}} &: \prod_{n,t,b_1,b_2}
    P(n,t) \to P(n,b_1) \to P(n,b_2) \to P(n, \mathsf{OCase}(t,b_1,b_2)) \\
  m_{\mathsf{Fix}} &: \prod_{n,b}
    P(n{+}1, b) \to P(n, \mathsf{OFix}(b))
\end{align}
(and similarly for $\mathsf{OSeq}$, $\mathsf{OLam}$, $\mathsf{OMap}$,
$\mathsf{OQuotient}$), there exists a unique function
\[
  \elim_{\mathsf{OTerm}} : \prod_{n:\Nat}\prod_{t:\mathsf{OTerm}(n)}
    P(n, t)
\]
computing by the expected case equations.
\end{definition}

\begin{proposition}[Pattern Matching = Elimination]
\label{prop:pattern-elim}
Every Python \texttt{if isinstance(t, OVar)} chain in
\texttt{denotation.py} is an instance of $\elim_{\mathsf{OTerm}}$.
The normalizer applies multiple elimination passes, each with a
different motive $P$:
\begin{enumerate}
  \item \textbf{Constant folding}: $P(n,t) = \mathsf{OTerm}(n)$,
    reduces $\mathsf{OOp}(\mathsf{add}, [\mathsf{OLit}(2),
    \mathsf{OLit}(3)])$ to $\mathsf{OLit}(5)$.
  \item \textbf{Algebraic simplification}: $P(n,t) = \mathsf{OTerm}(n)$,
    applies rewrite rules (e.g.\ $\mathsf{sorted} \circ \mathsf{sorted}
    \mapsto \mathsf{sorted}$).
  \item \textbf{Canonical form extraction}: $P(n,t) = \mathsf{String}$,
    produces the \texttt{canon()} hash for comparison.
  \item \textbf{Z3 compilation}: $P(n,t) = \mathsf{Z3Expr}$,
    translates OTerm to Z3 terms for the SMT solver.
\end{enumerate}
\end{proposition}

\section{Application: Invariant-Indexed Specifications}

Inductive families let us express specifications that depend on
structural properties of the input:

\begin{definition}[Sorted-Input Specification]
The type $\mathsf{Sorted}(A)$ of sorted lists:
\begin{align}
  \mathsf{snil} &: \mathsf{Sorted}(A) \\
  \mathsf{scons} &: \prod_{a : A} \mathsf{Sorted}(A) \to
    (\forall b \in \mathsf{head}.\; a \leq b) \to
    \mathsf{Sorted}(A)
\end{align}

A specification ``binary search on sorted input'' is:
\[
  S : \mathsf{Sorted}(\Int) \times \Int \to \mathsf{Option}(\Nat)
\]
The input is \emph{typed as sorted}, so the specification can assume
the precondition.
\end{definition}

\section{Automatic Precondition Inference via Inductive Families}

\begin{theorem}[Precondition Inference]
\label{thm:precondition}
Given a program $f$ and a specification $S$, if $f$ satisfies $S$
only on a subset of inputs, the \emph{precondition} $P$ can be
automatically inferred as:
\[
  P(x) = \mathsf{Spec}_S \text{ is satisfied at fiber } U_x
\]
where $U_x$ is the type fiber containing $x$.

The precondition is the set of fibers where the local \v{C}ech check
passes:
\[
  P = \{x : A \mid \Iso(\Sem_f, \Sem_S)(U_{\mathsf{tag}(x)}) = 1\}
\]
\end{theorem}

\begin{example}[Inferring ``requires positive integer'']
\begin{lstlisting}
def sqrt_int(n):
    if n < 0: return None  # error case
    r = 0
    while (r+1)*(r+1) <= n:
        r += 1
    return r
\end{lstlisting}

Specification: ``returns $\lfloor\sqrt{n}\rfloor$''.

On the $\mathsf{Int}^+$ fiber (positive integers): correct.
On the $\mathsf{Int}^-$ fiber (negative integers): returns None,
violating the spec.

Inferred precondition: $P(n) = (n \geq 0)$.

The cohomological decomposition reveals: $\Hone(U_{\mathsf{Int}^+}) = 0$
(correct on positives), $\Hone(U_{\mathsf{Int}^-}) \neq 0$
(fails on negatives).  The precondition is the union of passing fibers.
\end{example}

\section{Dependent Pattern Matching as Fiber Restriction}

Python's \texttt{isinstance} dispatch is a form of \emph{dependent
pattern matching}: the program branches on the type of its input,
and each branch has a different type signature.

In CTTT, this is \emph{exactly} the fiber decomposition:
\begin{itemize}
  \item Each \texttt{isinstance} branch restricts to a fiber.
  \item The branch body is the presheaf section on that fiber.
  \item The \v{C}ech complex checks that the branches are coherent
    (they agree on type boundaries).
\end{itemize}

\begin{theorem}[isinstance as Dependent Elimination]
\label{thm:isinstance-dep}
A function with \texttt{isinstance} dispatch:
\begin{lstlisting}
def f(x):
    if isinstance(x, int):    return h_int(x)
    elif isinstance(x, str):  return h_str(x)
    else:                     return h_other(x)
\end{lstlisting}
has the dependent type:
\[
  f : \prod_{x : \PyObj} B(\mathsf{tag}(x))
\]
where $B(\mathsf{TInt}) = \mathsf{codomain}(h_\mathsf{int})$,
$B(\mathsf{TStr}) = \mathsf{codomain}(h_\mathsf{str})$, etc.

The dependent type is inferred automatically from the \texttt{isinstance}
guards and the section extraction.
\end{theorem}

\begin{remark}[Dependent Pattern Matching in the Normalizer]
The normalizer in \texttt{denotation.py} itself uses dependent pattern
matching at \emph{two levels}:
\begin{enumerate}
  \item \textbf{On the input AST} (Python's \texttt{ast} module):
    each \texttt{ast.BinOp}, \texttt{ast.Call}, etc.\ maps to an
    OTerm constructor.  This is the compilation phase.
  \item \textbf{On the OTerm family}: each normalization pass
    matches on OTerm constructors and transforms them.  The arity
    index ensures that bound-variable operations (in $\mathsf{OFix}$
    and $\mathsf{OLam}$) correctly shift de Bruijn indices.
\end{enumerate}
The coherence of these two levels of pattern matching is guaranteed
by the $\v{C}$ech complex: the presheaf sections extracted from the
AST (level 1) are checked for global consistency via the coboundary
map on the OTerm normal forms (level 2).
\end{remark}


\chapter{The Computable Fragment of CTTT}
\label{ch:computable}

A critical question: which parts of CTTT are \emph{decidable}
(algorithms always terminate with a correct answer) vs.\
\emph{semi-decidable} (may not terminate) vs.\ \emph{undecidable}
(no algorithm exists)?

\section{The Three Decidability Classes}

\begin{definition}[Decidability Classification]
\label{def:decidability}
A problem $\Pi$ in the CTTT pipeline belongs to exactly one class:
\begin{enumerate}
  \item \textbf{Decidable.} There exists an algorithm that always
    terminates and returns the correct Yes/No answer.
  \item \textbf{Semi-decidable.} There exists an algorithm that
    terminates and returns Yes if the answer is positive, but may
    diverge (or return ``Unknown'') if the answer is negative.
  \item \textbf{Undecidable.} No algorithm can solve $\Pi$ for all
    inputs (by Rice's theorem or reduction from the halting problem).
\end{enumerate}
\end{definition}

\subsection{What is Decidable: Canonical Forms and Structural Equality}

The following operations are fully decidable because they operate on
finite, syntactic data:

\begin{itemize}
  \item \textbf{Duck type inference}: a single AST walk extracts the
    finite set of operations used on each parameter.  The duck-type
    lattice $\mathcal{D}$ is finite, so the meet/join is computed in
    $O(|A|)$ time.
  \item \textbf{Site category construction}: given $k$ parameters each
    with duck type sets of size $\leq d$, the site has at most $d^k$
    fibers.  Enumerating fibers and overlaps is $O(d^{2k})$.
  \item \textbf{Canonical form matching}: comparing two
    $\mathsf{OTerm}$ values via their \texttt{canon()} strings is
    decidable---it reduces to string equality on the normal forms.
  \item \textbf{$\Hone$ computation}: given the boolean fiber matrix,
    $\Hone$ is computed by GF(2) Gaussian elimination in $O(n^3)$.
\end{itemize}

\subsection{What is Semi-Decidable: Path Search}

\begin{proposition}[Path Search is Semi-Decidable]
\label{prop:path-semi}
The path search problem---``given $\den{f}$ and $\den{g}$, does
there exist a sequence of axiom applications connecting them?''---is
semi-decidable.  The BFS over the rewrite graph always terminates for
each depth bound $k$, but the search space may be infinite
(unboundedly many axiom applications), so the algorithm may not find
a path even though one exists at a depth not yet explored.
\end{proposition}

More precisely, the semi-decidable operations are:
\begin{itemize}
  \item \textbf{Local equivalence with folds}: Z3's
    $\mathsf{RecFunction}$ definitions may cause the SMT solver to
    diverge.  The fold RecFunction
    $\mathsf{fold}(\op, i, \mathit{stop}, \mathit{acc})$ recurses on
    $i$, and Z3 cannot always determine termination.
  \item \textbf{Local equivalence with axioms}: the 52 semantic axioms
    (Section~\ref{ch:path-space}) introduce universal quantifiers.
    Z3's E-matching engine may loop when instantiating
    $\forall x.\; \mathsf{sorted}(\mathsf{sorted}(x)) = \mathsf{sorted}(x)$
    on nested terms.
  \item \textbf{Bug finding per fiber}: satisfiability of the negated
    equivalence formula is in general undecidable for the full theory,
    though Z3 terminates on most practical instances.
\end{itemize}

\subsection{What is Undecidable: General Equivalence}

\begin{theorem}[Undecidability of General Equivalence]
\label{thm:undecidable-equiv}
The problem ``given arbitrary Python programs $f$ and $g$, are they
extensionally equivalent?''\ is undecidable.
\end{theorem}

\begin{proof}
By Rice's theorem.  Extensional equivalence is a non-trivial semantic
property of programs (it is not satisfied by all programs, nor
vacuously false).  Therefore no algorithm can decide it for all inputs.

Alternatively, reduce from the halting problem: given a Turing machine
$M$ and input $w$, let $f(x) = 1$ and $g(x) = [\text{run } M
\text{ on } w; \text{return } 1]$.  Then $f \equiv g$ iff $M$ halts
on $w$.
\end{proof}

CTTT does not attempt to solve the undecidable problem.  Instead, it
carves out a large \emph{decidable fragment} and returns ``Unknown''
for programs outside it, ensuring soundness without sacrificing
completeness on the fragment.

\section{Decidability Map}

\begin{center}
\small
\begin{longtable}{p{5cm} l p{5cm}}
  \toprule
  \textbf{Problem} & \textbf{Status} & \textbf{Justification} \\
  \midrule
  \endhead
  Duck type inference &
    Decidable &
    Finite AST walk, finite operation set \\
  Site category construction &
    Decidable &
    Finite product of finite tag sets \\
  Canonical form matching &
    Decidable &
    String equality on \texttt{canon()} output \\
  Structural OTerm equality &
    Decidable &
    Recursive descent on finite tree \\
  Local equivalence (QF) &
    Decidable &
    QF\_DT + QF\_LIA is decidable \\
  Local equivalence (with folds) &
    Semi-decidable &
    RecFunctions may cause non-termination \\
  Local equivalence (with axioms) &
    Semi-decidable &
    E-matching may loop \\
  Path search (depth-bounded) &
    Decidable &
    Finite BFS at each depth \\
  Path search (unbounded) &
    Semi-decidable &
    BFS terminates per level, infinite levels \\
  $\Hone$ computation &
    Decidable &
    GF(2) linear algebra, $O(n^3)$ \\
  Full pipeline (finite types) &
    Decidable &
    All components decidable \\
  Full pipeline (general) &
    Semi-decidable &
    Z3 may timeout on RecFunctions \\
  Equivalence (general) &
    Undecidable &
    Rice's theorem \\
  Spec checking (general) &
    Undecidable &
    Reduces to equivalence \\
  Abstract spec inference &
    Decidable &
    Finite abstract domains \\
  Bug finding (per fiber) &
    Semi-decidable &
    SAT check may timeout \\
  \bottomrule
\end{longtable}
\end{center}

\section{The Role of Z3}

Z3 provides the decision procedure for the first-order fragment of
CTTT.  Its role is precisely delineated:

\begin{definition}[The PyObj First-Order Theory]
\label{def:pyobj-theory}
The Z3 theory $\mathcal{T}_{\PyObj}$ is the first-order theory with:
\begin{itemize}
  \item \textbf{Sort}: $\PyObj$, defined as a Z3 algebraic datatype
    with seven constructors: $\mathsf{IntObj}(\mathit{ival} : \Int)$,
    $\mathsf{BoolObj}(\mathit{bval} : \Bool)$,
    $\mathsf{StrObj}(\mathit{sval} : \mathsf{String})$,
    $\mathsf{NoneObj}$, $\mathsf{Pair}(\mathit{fst}, \mathit{snd} : \PyObj)$,
    $\mathsf{Ref}(\mathit{addr} : \Int)$, $\mathsf{Bottom}$.
  \item \textbf{Tester predicates}: $\mathsf{is\_IntObj}$,
    $\mathsf{is\_BoolObj}$, \ldots---one per constructor.  These
    implement the $\mathsf{tag}$ fibration.
  \item \textbf{RecFunctions}: $\mathsf{binop}$, $\mathsf{unop}$,
    $\mathsf{truthy}$, $\mathsf{fold}$, $\mathsf{tag}$---each a
    recursively defined Z3 function.
  \item \textbf{Shared function symbols}: $\mathsf{py\_sorted}$,
    $\mathsf{py\_reversed}$, etc.---uninterpreted functions constrained
    by the 52 semantic axioms.
\end{itemize}
\end{definition}

\begin{proposition}[Decidability of the Quantifier-Free Fragment]
\label{prop:qf-decidable}
The quantifier-free fragment of $\mathcal{T}_{\PyObj}$ (without
RecFunction unfolding) is decidable.  It falls within the combined
theory QF\_DT + QF\_LIA + QF\_S (quantifier-free datatypes, linear
integer arithmetic, and strings), for which Z3 is a decision procedure.
\end{proposition}

\begin{proof}
By the Nelson-Oppen combination theorem.  QF\_DT is decidable
(finite case analysis on constructors), QF\_LIA is decidable
(Presburger arithmetic), and QF\_S is decidable (word equations
with length constraints).  The theories share only equality, so
their combination is decidable.
\end{proof}

\begin{remark}[RecFunctions as Controlled Undecidability]
The $\mathsf{fold}$ RecFunction introduces controlled undecidability:
Z3 unfolds the recursion up to a configurable depth bound.  When the
bound is reached without a proof, the solver returns ``Unknown''
rather than an incorrect answer.  This is sound (never produces
false equivalences) but incomplete (may miss valid equivalences that
require deeper unfolding).

In the implementation (\texttt{theory.py}), the fold is defined as:
\[
  \mathsf{fold}(\op, i, \mathit{stop}, \mathit{acc}) \coloneqq
  \begin{cases}
    \mathit{acc} & \text{if } i \geq \mathit{stop} \\
    \mathsf{binop}(\op, \mathsf{IntObj}(i),
      \mathsf{fold}(\op, i{+}1, \mathit{stop}, \mathit{acc}))
      & \text{otherwise}
  \end{cases}
\]
This is the computational content of the Nat-eliminator from
Chapter~\ref{ch:inductive-families}: $\mathsf{fold}$ is
$\rec_\Nat$ applied to the $\PyObj$ sort.
\end{remark}

\section{The Decidable Core}

The \emph{decidable core} of CTTT consists of programs where:
\begin{enumerate}
  \item All parameters have a single duck type (the spectral sequence
    degenerates---Theorem~\ref{thm:degeneration}).
  \item No recursive functions or unbounded loops (no RecFunctions
    appear in the compiled Z3 terms).
  \item No universal quantifiers in axioms (no E-matching required).
\end{enumerate}

For this core, the full pipeline (compilation, fiber decomposition,
local checks, $\Hone$ computation) is decidable with polynomial
complexity.

\begin{theorem}[Decidability of the Core]
\label{thm:decidable-core}
For programs in the decidable core, the equivalence/spec-checking/
bug-finding verdict is computed in time:
\[
  O(|A_f| + |A_g| + |\site|^3)
\]
where $|A_f|, |A_g|$ are the AST sizes and $|\site|$ is the number
of sites (type fiber combinations).
\end{theorem}

\begin{proof}
\begin{enumerate}
  \item AST compilation to OTerm: $O(|A|)$ (single AST walk,
    each node maps to one OTerm constructor via the compilation
    function in \texttt{denotation.py}).
  \item Duck type inference: $O(|A|)$ (single walk, collecting
    operations into the finite lattice $\mathcal{D}$).
  \item Site construction: $O(|\site|^2)$ (enumerate products and
    overlaps of duck-type fibers).
  \item Local checks: $O(|\site|)$ Z3 queries, each in QF\_DT +
    QF\_LIA (decidable by Proposition~\ref{prop:qf-decidable},
    polynomial with fixed formula size).
  \item $\Hone$ computation: $O(|\site|^3)$ (GF(2) Gaussian
    elimination on the coboundary matrix).
\end{enumerate}
Total: dominated by the $\Hone$ computation, which is cubic in the
number of fibers.
\end{proof}

\section{Extending the Decidable Fragment}

Each dichotomy theory (Part V) extends the decidable fragment by
adding axioms that convert semi-decidable checks to decidable ones:

\begin{itemize}
  \item \textbf{D1 (Rec$\leftrightarrow$Iter)}: The Nat-eliminator
    extraction converts RecFunctions to canonical folds, making the
    local check decidable (syntactic equality of fold terms).
  \item \textbf{D2 (Nested$\leftrightarrow$Flat)}: Function inlining
    eliminates nested calls, producing ground terms.
  \item \textbf{D4 (Comp$\leftrightarrow$Loop)}: Comprehension hashing
    reduces to shared symbol equality (decidable).
  \item \textbf{D8 (Multi-Return$\leftrightarrow$Single)}: Section
    merging reduces to Z3 If-expression comparison (decidable).
\end{itemize}

Each extension enlarges the set of programs for which the pipeline
is fully decidable, without compromising soundness for programs
outside the fragment (they get ``Unknown'' rather than wrong answers).

\begin{proposition}[Monotonic Decidability Growth]
\label{prop:monotonic-decidability}
Let $\mathcal{C}_0$ be the decidable core and $\mathcal{C}_k$ the
fragment after adding dichotomy theories $D_1, \ldots, D_k$.  Then:
\[
  \mathcal{C}_0 \subseteq \mathcal{C}_1 \subseteq \cdots
    \subseteq \mathcal{C}_k
\]
and each inclusion is strict (each dichotomy theory brings at least
one previously semi-decidable program class into the decidable
fragment).  Furthermore, the decidable core is \emph{never}
reduced: adding axioms can only enlarge the fragment because each
axiom is a sound equational law.
\end{proposition}


